\subsection{SPACES OF PROJECTIVE LINEAR VARIETIES}

Given a division ring K, the set \(\mathbf{P}_{n,p}(\mathbf{K})\) of projective linear varieties of \(p \geqslant 0\) dimensions of the left projective space \(\mathbf{P}_{n}(\mathbf{K})\) is clearly in one-to-one correspondence with the set of vector subspaces of \(p + 1\) dimensions of the left vector space \(K_{s}^{n+1}\) . Let \(L_{n+1, p+1}(\mathbf{K})\) denote the set of free systems \((x_{k})_{1 \leqslant k \leqslant p+1}\) of \(p + 1\) vectors of \(K_{s}^{n+1}\) ; then the set \(\mathbf{P}_{n,p}(\mathbf{K})\)

is again in one-to-one correspondence with the quotient of \(\mathbf{L}_{n+1,p+1}(\mathbf{K})\) by the equivalence relation \(\Delta_{n,p}(\mathbf{K})\) : `` \((\mathbf{x}^{k})\) and \((\mathbf{y}^{k})\) generate the same vector subspace of \(p+1\) dimensions of \(K_{s}^{n+1}\) ''. In what follows we shall identify \(\mathbf{P}_{n,p}(\mathbf{K})\) with this quotient set. On the other hand, if to each free system \((\mathbf{x}_{k})\) of \(p+1\) vectors of \(K_{s}^{n+1}\) we make correspond the matrix X of \(p+1\) rows and \(n+1\) columns for which \(x_{k}\) is the kth row ( \(1\leqslant k\leqslant p+1\) ), we have a one-to-one correspondence between \(L_{n+1,p+1}(\mathbf{K})\) and the set of all matrices of \(p+1\) rows and \(n+1\) columns which are of rank \(p+1\) ; we shall identify \(L_{n+1,p+1}(\mathbf{K})\) with this set of matrices, and the relation \(\Delta_{n,p}(\mathbf{K})\) between two matrices X, Y is then the following: `` there exists a non-singular square matrix T of order \(p+1\) such that Y = T.X''.

In what follows we shall take K to be the field R, and we shall omit the letter K in the notation above. We can define a topology on \(P_{n,p}\) by a process which generalizes the definition of the topology of real projective spaces. Namely, \(L_{n+1, p+1}\) is contained in the space \(M_{p+1, n+1}\) of all matrices of \(p + 1\) rows and \(n + 1\) columns with real elements; we endow \(L_{n+1, p+1}\) with the topology induced by the topology of this matrix space (§ 1, no. 6).

DEFINITION 2. The space \(\mathbf{P}_{n,p}\) which is the quotient of the topological space \(\mathbf{L}_{n+1,p+1}\) by the equivalence relation \(\Delta_{n,p}\) is called the space of projective linear varieties of \(p\geqslant0\) dimensions in real projective space \(\mathbf{P}_{n}\) .

We shall use the following notation: given a matrix X of \(p + 1\) rows and \(n + 1\) columns, and any strictly increasing sequence

\[
\sigma = (i _ {1}, \dots , i _ {p + 1})
\]

of \(p + 1\) indices belonging to the interval \([0, n]\) of \(\mathbf{N}\) , we denote by \(X_{\sigma}\) the square submatrix of \(X\) formed by the columns whose indices are \(i_1, i_2, \ldots, i_{p+1}\) . We denote by \(\mathbf{A}_{\sigma}\) the subset of \(\mathbf{L}_{n+1, p+1}\) consisting of these matrices \(X\) such that \(X_{\sigma}\) is non-singular. By Proposition 6 of § 1, no. 6, \(\mathbf{A}_{\sigma}\) is a dense open set in \(\mathbf{M}_{p+1, n+1}\) , and the function \(X \to X_{\sigma}^{-1}\) is continuous on \(\mathbf{A}_{\sigma}\) .

A geometrical interpretation of the set \(A_{\sigma}\) is as follows: let \(E_{\sigma}\) be the vector subspace of \(R^{n+1}\) generated by the vectors \(e_{i}\) of the canonical basis such that \(i \in \sigma\) , and let \(E_{\sigma}'\) be the complementary subspace generated by the \(e_{i}\) such that \(i \notin \sigma\) ; then to say that a matrix x belongs to \(A_{\sigma}\) means that the projections on \(E_{\sigma}\) of its \(p + 1\) rows of \(x_{k}\) form a free system, or again that the vector subspace generated by the \(x_{k}\) is a complement of \(E_{\sigma}'\) (or that its intersection with \(E_{\sigma}'\) consists only of o).

PROPOSITION 6. The space \(\mathbf{P}_{n,p}\) is Hausdorff.

We show first that the relation \(\Delta_{n,p}\) is open. If U is an open set in \(L_{n+1,p+1}\) , to saturate U with respect to \(\Delta_{n,p}\) we have to take the union of the images of U under the mapping \(X \to T \cdot X\) , where T runs through the set of non-singular square matrices of order \(p + 1\) ; since each of these mappings is bicontinuous, all these images are open sets and therefore so is their union.

By Proposition 8 of Chapter I, § 8, no. 3, the proof will be complete if we show that the subset N of \(L_{n+1,p+1} \times L_{n+1,p+1}\) , defined by \(\Delta_{n,p}\) , is closed. Let \((X,\mathcal{Y})\) be a point of the product space which lies in the closure of N, and let \(\sigma\) be a sequence of indices such that \(X_{\sigma}\) is non-singular: since \(A_{\sigma}\) is open, there is a neighbourhood V of \((X,\mathcal{Y})\) such that, for each pair \((X',\mathcal{Y}') \in \mathrm{N} \cap \mathrm{V}\) , the matrix \(X'_{\sigma}\) is non-singular; as \((X',\mathcal{Y}')\) tends to \((X,\mathcal{Y})\) while remaining in N, the matrix \(Y'_{\sigma}X_{\sigma}^{\prime -1}\) therefore tends to \(T = Y_{\sigma}X_{\sigma}^{-1}\) ; since we have \(Y' = (Y'_{\sigma}X_{\sigma}^{\prime -1})X'\) , we see by passing to the limit that Y = T.X, and the proof is complete.

PROPOSITION 7. The space \(\mathbf{P}_{n,p}\) is compact.

It is enough to show that there is a compact subspace of \(L_{n+1, p+1}\) which meets every equivalence class mod \(\Delta_{n,p}\) in at least one point; for \(P_{n,p}\) is then the image of this subspace under the canonical mapping of \(L_{n+1, p+1}\) onto \(P_{n,p}\) , and is therefore compact (Chapter I, § 9, no. 4, Theorem 2).

Let \(\mathbf{V}_{n+1, p+1}\) be the subspace of \(\mathbf{L}_{n+1, p+1}\) whose elements are the systems \((\boldsymbol{x}_k)\) of \(p + 1\) vectors forming an orthonormal Euclidean basis of the vector subspace they generate that is to say such that \((x_h | x_k) = 0\) whenever \(h \neq k\) , \((x_h | x_h) = 1\) for \(1 \leqslant h \leqslant p + 1\) . Every vector subspace of \(p + 1\) dimensions of \(\mathbf{R}^{n+1}\) has such a basis and therefore every class mod. \(\Delta_{n,p}\) meets \(\mathbf{V}_{n+1, p+1}\) . On the other hand, the matrices \(X = (x_{ij})\) of \(\mathbf{V}_{n+1, p+1}\) are defined by the relations

\[
\sum_ {j = 0} ^ {n} x _ {i j} ^ {2} = 1 (1 \leqslant i \leqslant p + 1),
\]

\[
\sum_ {j = 0} ^ {n} x _ {i j} x _ {k j} = 0 \quad (i \neq k);
\]

therefore they form a closed set in \(M_{p+1, n+1}\) ; and since these relations imply that \(|x_{ij}| \leqslant 1\) for each pair of indices \((i, j)\) , this set is bounded and hence compact.

PROPOSITION 8. \(\mathbf{P}_{n,p}\) is connected and locally connected, and every point has an open neighbourhood homeomorphic to \(\mathbf{R}^{(p+1)(n-p)}\) .

For every (strictly increasing) sequence of indices \(\sigma\) , the set \(\mathbf{A}_{\sigma}\) is open in \(\mathbf{L}_{n+1, p+1}\) and saturated with respect to \(\Delta_{n,p}\) ; its canonical image

\(C_{\sigma}\) in \(P_{n,p}\) is therefore an open set homeomorphic to the quotient of \(A_{\sigma}\) by the equivalence relation \(\Theta_{\sigma}\) induced on \(A_{\sigma}\) by \(\Delta_{n,p}\) (Lemma 1).

Let \(\mathbf{B}_{\sigma}\) be the subset of \(\mathbf{A}_{\sigma}\) consisting of matrices \(X\) such that \(X_{\sigma}\) is the unit matrix of order \(p + 1\) ; the elements of \(X\) other than those of \(X_{\sigma}\) are then arbitrary, and therefore \(\mathbf{B}_{\sigma}\) is homeomorphic to the space \(\mathbf{R}^{(p + 1)(n - p)}\) . To each matrix \(X \in \mathbf{A}_{\sigma}\) let us make correspond the matrix \(\Upsilon = X_{\sigma}^{-1}X\) , which belongs to \(\mathbf{B}_{\sigma}\) ; then we have defined a continuous mapping \(g\) of \(\mathbf{A}_{\sigma}\) onto \(\mathbf{B}_{\sigma}\) , such that \(g(X)\) is the only matrix of \(\mathbf{B}_{\sigma}\) congruent to \(X\) mod \(\Theta_{\sigma}\) . It follows that \(\mathbf{B}_{\sigma}\) is homeomorphic to \(\mathbf{A}_{\sigma}/\Theta_{\sigma}\) (Lemma 2), hence to \(\mathbf{C}_{\sigma}\) .

The set \(C_{\sigma}\) is therefore connected. Since \(A_{\sigma}\) is dense in \(L_{n+1, p+1}\) , \(C_{\sigma}\) is dense in \(P_{n, p}\) and therefore \(P_{n, p}\) is also connected (Chapter I, § 11, no. 1, Proposition 1). On the other hand, every point of \(P_{n, p}\) belongs to \(C_{\sigma}\) for at least one sequence of indices \(\sigma\) , and therefore has an open neighbourhood homeomorphic to \(\mathbf{R}^{(p+1)(n-p)}\) .

The matrix \(Y=g(X)\) can be interpreted as follows: let us suppose for the sake of simplicity that the sequence \(\sigma\) consists of the \(p+1\) indices n-p, \(n-p+1\) , \(\ldots\) , n, and let \(a_{ij}\) ( \(1\leqslant i\leqslant p+1\) , \(o\leqslant j\leqslant n-p-1\) ) denote the elements of the first n-p columns of y; then the vector subspace of \(R^{n+1}\) generated by the rows of x is that defined by the equations

\[
x _ {j} = \sum_ {i = 1} ^ {p + 1} a _ {i j} x _ {n - p + i - 1} \quad (0 \leqslant j \leqslant n - p - 1).
\]

